\section{Experiment 6: NeuroGolf Trace Language and MAS Verification}
\label{sec:experiment-6}

We apply the trace-language framework to a \emph{multi-agent system (MAS)} that
builds ONNX-formulated neural-network solvers for the Kaggle NeuroGolf 2026
competition\footnote{\url{https://kaggle.com/competitions/neurogolf-2026}}.
ARC-AGI style grid-transformation tasks are submitted as ONNX \texttt{.tar.gz}
bundles whose score is determined by the cost function
\(\text{points} = \max(1, 25 - \ln(\max(1, \text{params} + \text{memory})))\).

\subsection{Operation Alphabet \(\Sigma\)}
\label{sec:exp6-alphabet}

Table~\ref{tab:exp6-alphabet} defines the 20-operation alphabet for the
NeuroGolf domain. Operations span five phases: analysis~(\textsc{analyze},
\textsc{discover}), construction~(\textsc{build}, \textsc{encode},
\textsc{convolution}, \textsc{label\_propagate}, \textsc{scatternd}),
optimization~(\textsc{fp16\_surgery}, \textsc{reduce\_fusion},
\textsc{cast\_collapse}, \textsc{dtype\_narrow}, \textsc{prune}),
verification and costing~(\textsc{verify\_train}, \textsc{verify\_test},
\textsc{verify\_arc\_gen}, \textsc{compute\_cost}, \textsc{cost\_grader\_match}),
and submission~(\textsc{blend\_bundle}, \textsc{package},
\textsc{submit}, \textsc{reject}, \textsc{halt}).

\begin{table}[htbp]
\centering
\caption{NeuroGolf operation alphabet.}
\label{tab:exp6-alphabet}
\small
\begin{tabular}{lll}
\toprule
Operation & Agent Role & Description \\
\midrule
\textsc{analyze\_task}    & Task Analyzer   & Inspect grid I/O pairs for structure \\
\textsc{discover\_pattern}& Pattern Miner   & Match ARC-AGI transformation catalog \\
\textsc{build\_onnx}      & ONNX Builder    & Construct hand-built ONNX graph \\
\textsc{encode\_rule}     & ONNX Builder    & Encode transformation as ONNX ops \\
\textsc{convolution}       & ONNX Builder    & Apply convolution-based solver \\
\textsc{label\_propagate}  & ONNX Builder    & Label-propagation via MaxPool \\
\textsc{scatternd\_hist}   & ONNX Builder    & ScatterND-based histogram count \\
\textsc{fp16\_surgery}    & Optimizer       & Convert FLOAT tensors to FP16 \\
\textsc{reduce\_fusion}   & Optimizer       & Fuse ReduceSum chains \\
\textsc{cast\_collapse}   & Optimizer       & Collapse adjacent Cast nodes \\
\textsc{dtype\_narrow}    & Optimizer       & Narrow bool/int tensors \\
\textsc{prune}            & Optimizer       & Remove dead nodes \\
\textsc{verify\_train}    & Verifier        & Check correctness on train set \\
\textsc{verify\_test}     & Verifier        & Check correctness on test set \\
\textsc{verify\_arc\_gen} & Verifier        & Check correctness on ARC-GEN set \\
\textsc{compute\_cost}    & Cost Grader     & Compute params + memory cost \\
\textsc{cost\_grader\_match} & Cost Grader  & Compare cost to grader formula \\
\textsc{blend\_bundle}    & Blender         & Select cheapest ONNX per task \\
\textsc{package}          & Packager        & Create \texttt{submission.tar.gz} \\
\textsc{submit}           & Orchestrator    & Submit bundle to Kaggle \\
\textsc{reject}           & Any             & Reject unsolvable task \\
\textsc{halt}             & Any             & Terminate pipeline \\
\bottomrule
\end{tabular}
\end{table}

\subsection{DFA Verifier}
\label{sec:exp6-dfa}

The verifier is a DFA \(V = (Q, \Sigma, \delta, q_0, F)\) where
\(|Q| = 13\) states capture the phases of the pipeline:
\(\{\textsc{init}, \textsc{analyzed}, \textsc{built}, \textsc{optimized},
\textsc{verified\_train}, \textsc{verified\_test}, \textsc{verified\_arc\_gen},
\textsc{costed}, \textsc{blended}, \textsc{packaged}, \textsc{submitted},
\textsc{rejected}, \textsc{error}\}\).
The acceptance states are \(F = \{\textsc{submitted}, \textsc{packaged}\}\).

Critical constraints enforced by \(\delta\) include:
\begin{enumerate}
  \item An \textsc{analyze\_task} must precede any construction operation.
  \item \textsc{verify\_train} \(\rightarrow\) \textsc{verify\_test} \(\rightarrow\)
        \textsc{verify\_arc\_gen} are ordered.
  \item A \textsc{reject} transitions directly to the terminal rejecting state.
  \item \textsc{submit} can only follow \textsc{package}.
  \item \textsc{blend\_bundle} requires at least \textsc{verify\_test} completion.
\end{enumerate}

\subsection{Scenarios and Results}
\label{sec:exp6-scenarios}

Table~\ref{tab:exp6-scenarios} summarizes five verification scenarios.

\begin{table}[htbp]
\centering
\caption{Verification results for five NeuroGolf MAS scenarios.}
\label{tab:exp6-scenarios}
\small
\begin{tabular}{lcc}
\toprule
Scenario & Final State & Accepted \\
\midrule
  Correct Pipeline (13 tasks)              & SUBMITTED            & \checkmark \\
  Build w/o Analysis                       & ERROR                & \ding{55} \\
  Reject Unsolvable Task                   & REJECTED             & \ding{55} \\
  Skip ARC-GEN Verification                & SUBMITTED            & \ding{55} \\
  Blend Before Full Verify                 & ERROR                & \ding{55} \\
\bottomrule
\end{tabular}
\end{table}

\subsection{Keyword Coverage Analysis}
\label{sec:exp6-keywords}

We scanned the top-4 Kaggle notebooks for 11 methodology keywords.
Table~\ref{tab:exp6-keywords} reports the coverage per notebook.

\begin{table}[htbp]
\centering
\caption{Keyword coverage ratios across top NeuroGolf notebooks.}
\label{tab:exp6-keywords}
\small
\begin{tabular}{lcc}
\toprule
Notebook & Keywords Found & Coverage Ratio \\
\midrule
\bottomrule
\end{tabular}
\end{table}

\subsection{Discussion}
\label{sec:exp6-discussion}

Experiment~6 demonstrates that the trace-language framework generalizes
beyond the toy implementations of Experiments~1--5 to a realistic,
competitive multi-agent system. The DFA verifier successfully distinguishes
correct pipelines from malformed traces (missing analysis, skipped
verification, premature blending).

The keyword coverage analysis reveals which methodology patterns the top
entries prioritize: notebooks emphasizing ONNX structure (convolution,
label propagation, scatternd) and cost optimizations (FP16 surgery, reduce
fusion, cast collapse) achieve higher coverage. Missing keywords typically
correspond to unused techniques (e.g., ARC-GEN-specific verification or
certain pruning strategies).

A limitation is that keyword coverage is a surface-level proxy: the presence
of a keyword in notebook code does not guarantee that the technique is
implemented correctly, and absence does not preclude equivalent implicit
implementations. Nevertheless, the combination of DFA-based operation
verification and keyword coverage provides a multi-faceted view of MAS
completeness.
